\documentclass[11pt]{article}
\usepackage[margin=1.15in]{geometry}
\usepackage{amsmath,amssymb,graphicx,xcolor,needspace}
\usepackage[hidelinks]{hyperref}
\newcommand{\Cov}{\mathsf C}
\newcommand{\T}{^{\mathsf T}}
\newcommand{\us}{\_\allowbreak}
\emergencystretch=3em
\begin{document}
\begin{center}
{\LARGE Baryonic feedback in a joint shear forecast\par}
\vspace{.4em}
{\large Power, bispectrum, neutrinos and the information added by Map3\par}
\vspace{.6em}
{\large November 2026\par}
\end{center}

\section{The question and the companion notebook}
\paragraph{Moving matter changes lensing.}
Gas can cool into galaxies or be driven out of halos by energetic stellar
and black-hole activity. These processes redistribute matter. A lensing
prediction based on collisionless gravitational evolution alone may then
misestimate small-scale correlations. Both the power spectrum and the
bispectrum respond, but their responses need not be related by a single
amplitude. A joint shear analysis can therefore test feedback parameters
through two different patterns in the data.

\paragraph{What this study attempts.}
The methods paper by Gomes and collaborators~\cite{methods} tested scale
cuts against a specified baryonic contamination. Our notebook asks a
related but different question: if a paired response model is included
explicitly, does the aperture statistic improve local constraints on its
feedback coordinates? Burger and collaborators~\cite{burger} provide the
paired emulation used here. This is not a claim that a five-parameter
Fisher calculation reproduces either paper's posterior or establishes the
accuracy of the emulator outside its training region.

\paragraph{One notebook contains the analysis.}
Open
\href{EXAMPLE_FISHER_BACCO.ipynb}{\texttt{EXAMPLE\us FISHER\us BACCO.ipynb}}
beside this note in the activated Cocoa environment. It reads the numbered
TATT example, enables the additional BACCO theory block, and calls Cobaya
for every prediction. It then uses the shared Fisher functions and plots
the means, derivative responses and parameter errors. The original
absolute-bispectrum provider remains in the YAML. Expensive cells run real
models; their outputs are empty following the source refactor. No new
physical forecast is reported in this note. Generated results belong in
\texttt{chains/fisher\us bacco/}, not among the scientific input files.

\section{Why the joint calculation needs two response functions}
\paragraph{Separate an absolute spectrum from its response.}
Let superscript $D$ denote the gravity-only reference and $H$ the prediction
including the chosen gas physics. The power response is defined as
\begin{equation}
 R_P(k,z)=\frac{P^H(k,z)}{P^D(k,z)}.
\end{equation}
It is dimensionless. Once an absolute reference spectrum is supplied,
recover the modified spectrum by multiplication:
\begin{equation}
 P^H(k,z)=R_P(k,z)P^D(k,z).
\end{equation}
A response alone does not specify the absolute spectrum. Different
reference cosmologies can have the same response and different powers.

\paragraph{The bispectrum needs triangle information.}
A bispectrum depends on three Fourier wavevectors whose sum is zero.
Their lengths form a triangle. Its response is
\begin{equation}
 R_B(k_1,k_2,k_3,z)=\frac{B^H(k_1,k_2,k_3,z)}{B^D(k_1,k_2,k_3,z)}.
\end{equation}
The corresponding absolute prediction is
\begin{equation}
 B^H(k_1,k_2,k_3,z)=R_B(k_1,k_2,k_3,z)B^D(k_1,k_2,k_3,z).
\end{equation}
The two-point shear projection uses the power spectrum, while three-point
shear uses the bispectrum. Supplying only $R_B$ would leave the feedback
model inconsistent between the two observables. Equally, assuming
$R_B=R_P^{3/2}$ would impose a relation that an independently emulated
bispectrum can test rather than assume.

\paragraph{The theory blocks have different jobs.}
CAMB supplies the cosmology and spectra requested by the model. The
absolute-bispectrum block supplies $B^D$. The separate
\texttt{bacco\us bispectrum} block supplies the paired responses. The
likelihood requests those products and projects the resulting model.
BiHalofit is the supplied absolute-bispectrum example, not a mandatory
identity of the three-point likelihood: another provider can satisfy the
same physical input requirement. Enabling BACCO does not justify applying
its power ratio a second time in an independent feedback block.

\paragraph{What a feedback coordinate means.}
The notebook varies logarithmic mass and gas-ejection coordinates:
\begin{equation}
 p_M=\log_{10}\left[\frac{M_c}{h^{-1}M_\odot}\right].
\end{equation}
The mass scale $M_c$ controls gas retention in the model. The other varied
coordinate is
\begin{equation}
 p_\eta=\log_{10}\eta.
\end{equation}
The dimensionless parameter $\eta$ describes gas ejection. A change of
$0.1$ in a base-ten logarithm multiplies the original quantity by
$10^{0.1}\simeq1.26$; it does not add $0.1$ to that quantity. The other
feedback coordinates remain fixed at their declared values in this
five-parameter experiment.

\section{Cold matter and total matter are not interchangeable}
\paragraph{Start with the density fields.}
Here ``cold'' includes cold dark matter and baryons. Massive neutrinos
contribute to the total matter seen by lensing as well. Define the mass
fractions
\begin{equation}
 f_\nu=\frac{\Omega_\nu}{\Omega_m}.
\end{equation}
The remaining fraction is
\begin{equation}
 f_c=1-f_\nu.
\end{equation}
The total density contrast is the weighted sum
\begin{equation}
 \delta_m=f_c\delta_c+f_\nu\delta_\nu.
\end{equation}
These weights are background mass fractions. They are not the fractions
of a particular Fourier mode's fluctuation amplitude.

\paragraph{Expand the power before choosing an approximation.}
Multiplying two copies of this field produces four terms. The two mixed
terms have the same power, so they combine to give
\begin{equation}
 P_m^D=f_c^2P_c^D+2f_cf_\nu P_{c\nu}^D+f_\nu^2P_{\nu\nu}^D.
\end{equation}
The default implementation uses CAMB's linear massive-neutrino power and
approximates the cross-spectrum by aligned fluctuation phases:
\begin{equation}
 P_{c\nu}^D\simeq\sqrt{P_c^D P_{\nu\nu,L}}.
\end{equation}
Substituting this expression turns the three terms into a square:
\begin{equation}
 P_m^D\simeq
 \left(f_c\sqrt{P_c^D}+f_\nu\sqrt{P_{\nu\nu,L}}\right)^2.
\end{equation}
This is an approximation to the cross-correlation, not an identity obeyed
by arbitrary nonlinear cold and neutrino fields.

\paragraph{Preserve the supplied gravity-only spectrum.}
The code starts from the existing total nonlinear power rather than
replacing it by a second background calculation. Define its amplitude
\begin{equation}
 U=\sqrt{P_m^D}.
\end{equation}
Define the mass-weighted neutrino amplitude from the explicit massive-neutrino
spectrum:
\begin{equation}
 N=f_\nu\sqrt{P_{\nu\nu,L}}.
\end{equation}
The inferred cold amplitude is then
\begin{equation}
 U_c=U-N.
\end{equation}
All three amplitudes have units of the square root of a power spectrum.
The construction requires $0\leq N\leq U$. Silently replacing a negative
inferred cold amplitude by zero would change this model. In the absence
of massive neutrinos we set $N=0$ directly.

\paragraph{Apply the cold response at the amplitude level.}
BACCO supplies the cold power ratio. Its amplitude ratio is
\begin{equation}
 q=\sqrt{R_{P,c}}.
\end{equation}
Holding the neutrino auto-spectrum fixed gives the total hydro amplitude
$qU_c+N$. Therefore
\begin{equation}
 R_{P,m}=\frac{(qU_c+N)^2}{U^2}.
\end{equation}
To see the relative weights more clearly, define
\begin{equation}
 a=U_c/U.
\end{equation}
The neutrino fluctuation fraction is
\begin{equation}
 n=N/U=1-a.
\end{equation}
These quantities depend on wavenumber and redshift; $a$ here is not the
cosmic scale factor. Substituting gives
\begin{equation}
 R_{P,m}=(aq+n)^2.
\end{equation}

\paragraph{Work one example through.}
Choose $a=0.99$ and $R_{P,c}=0.81$ solely to illustrate the arithmetic.
Then $q=0.9$ and $n=0.01$. The total amplitude ratio is
\begin{equation}
 aq+n=0.99\times0.9+0.01=0.901.
\end{equation}
Squaring it gives
\begin{equation}
 R_{P,m}=0.901^2=0.811801.
\end{equation}
The unmodified neutrino contribution slightly dilutes the cold suppression.
The alternative \texttt{cold\us on\us total} prescription instead sets
$R_{P,m}=R_{P,c}$. The notebook compares these two assumptions with the same
reference cosmology. The numerical values in this paragraph are not an
emulator evaluation.

\section{Why the bispectrum conversion is an additional assumption}
\paragraph{There are eight ordered field contributions.}
Each of three density factors can be cold or neutrino. Expanding their
product gives
\begin{equation}
\begin{split}
 B_m={}&f_c^3B_{ccc}
 +f_c^2f_\nu(B_{cc\nu}+B_{c\nu c}+B_{\nu cc})\\
 &+f_cf_\nu^2(B_{c\nu\nu}+B_{\nu c\nu}+B_{\nu\nu c})
 +f_\nu^3B_{\nu\nu\nu}.
\end{split}
\end{equation}
The position of each subscript identifies its Fourier leg. Thus
$B_{c\nu c}$ places the neutrino field on the second side. For unequal
sides, its value need not equal the term with the neutrino on the first
side. BACCO's cold response alone does not determine all these terms.

\paragraph{The implemented assignment.}
For each leg compute its own $a_i,n_i,q_i$ from the power prescription.
If every cold leg simply acquired its power-amplitude response, the
three-leg response would be
\begin{equation}
 Q=\prod_{i=1}^3(a_iq_i+n_i).
\end{equation}
This product contains eight terms. Its pure-cold term is
$a_1a_2a_3q_1q_2q_3$. Remove that term and replace it by the independently
emulated cold bispectrum ratio. The result is
\begin{equation}
 R_{B,m}=Q-a_1a_2a_3q_1q_2q_3+a_1a_2a_3R_{B,c}.
\end{equation}
Combining the last two terms gives the expression used for checking:
\begin{equation}
 R_{B,m}=\prod_{i=1}^3(a_iq_i+n_i)
 +\left(\prod_{i=1}^3a_i\right)
  \left(R_{B,c}-\prod_{i=1}^3q_i\right).
\end{equation}

\paragraph{Checks of the formula are not a physical calibration.}
If every response equals one, each first-product factor is $a_i+n_i=1$
and the second term vanishes. Thus $R_{B,m}=1$ and the supplied baseline is
restored. If neutrinos are absent, all $a_i=1$ and all $n_i=0$. The power
products cancel and the answer is exactly $R_{B,c}$. These checks verify
limiting cases. They do not determine the actual mixed cold-neutrino
bispectra. In particular, assigning a term proportional to
$n_1n_2n_3B_m^D$ is not the bispectrum of a strictly Gaussian linear
neutrino field, whose bispectrum would vanish. ``Linear neutrinos'' in
this implementation names the power input, not a derivation of every
three-point contribution.

\paragraph{Which fluctuation amplitude enters the emulator.}
The forecast's $\sigma_8$ is the total linear-density RMS in spheres of
radius $8\,h^{-1}\mathrm{Mpc}$. BACCO's cold coordinate instead uses the
cold linear spectrum. CAMB must supply that coordinate at each displaced
cosmology. Holding the cold coordinate fixed while varying total $\sigma_8$
would miss part of the feedback response. The current emulator also has
no independent coordinate for every cosmological parameter. A changed
CAMB prediction does not by itself establish that the feedback model's
cosmology dependence is complete.

\section{Continuing beyond training and testing numerical sampling}
\paragraph{Triangle shape must be specified during continuation.}
Sort the side lengths from shortest to longest. One shape coordinate is
\begin{equation}
 u=k_{\rm short}/k_{\rm middle}.
\end{equation}
A second coordinate is
\begin{equation}
 v=(k_{\rm long}-k_{\rm middle})/k_{\rm short}.
\end{equation}
It is a side-ratio coordinate, not an opening angle. Scaling all sides by
the same positive factor preserves both $u$ and $v$. Clamping each side
separately to a boundary generally changes them. For example, the triangle
$(2,6,7)$ scaled to longest side five becomes $(10/7,30/7,5)$; independently
clamping its long sides at five gives $(2,5,5)$. They are different shapes.

\paragraph{What the notebook compares.}
The physical comparisons change the neutrino prescription, the high-$k$
continuation, its anchor, the low-$k$ anchor, or the high-redshift
continuation. These choices assign predictions outside measured training
support. A denser numerical table cannot establish which continuation is
physically correct. The numerical comparisons instead refine the redshift,
wavenumber and two shape grids, or the dense interpolation factor, while
holding the selected continuation fixed. The notebook labels which cases
compute derivatives and which compare only means; it does not infer
full derivative convergence from a small mean difference.

\paragraph{Measure a difference with the survey covariance.}
For a changed mean define
\begin{equation}
 \Delta=\mu_{\rm changed}-\mu_{\rm reference}.
\end{equation}
The full squared distance is
\begin{equation}
 \Delta\chi^2=\Delta\T C^{-1}\Delta.
\end{equation}
Here $C$ is the fixed joint covariance. Retain its cross block. The
companion intrinsic-alignment note derives how this distance separates
into a two-point contribution and the additional conditional three-point
contribution. A large distance identifies model sensitivity, not a
parameter bias after fitting. Estimating a fitted bias would require
projecting the difference onto the allowed parameter responses.

\section{What the five-parameter forecast means}
\paragraph{State the restricted parameter space first.}
The varied coordinates are total $\sigma_8$, $p_M$, $p_\eta$, $A_1$ and
$A_2$. The remaining cosmological, alignment, source and gas parameters
are fixed. There are no added Gaussian priors on these five coordinates.
The full five-by-five Fisher inverse marginalizes over the other four
varied coordinates when one error is read; it remains conditional on all
coordinates held fixed outside this list. Comparing such errors with a
forecast that releases every nuisance would compare different experiments.

\paragraph{Derivatives and displayed responses have different steps.}
The five-point derivatives use small absolute parameter displacements.
The display uses larger chosen changes, such as $0.2$ in $p_M$ and $0.1$
in $p_\eta$, to show derivative patterns in row-error units. Multiplying
a derivative by a display displacement is a tangent prediction. It is
not a separately evaluated finite-offset model. The best-fit response
coefficient uses the entire covariance even when a figure shows only one
source-bin combination.

\paragraph{Parameter errors can reach an emulator boundary.}
The Fisher approximation replaces the local likelihood by a quadratic.
Its Gaussian tails can extend beyond a permitted emulator coordinate
range. The notebook compares the inferred interval with the declared
coordinate bounds. Those bounds define this example's parameter box;
checking them does not independently validate every emulator training
coordinate. A crossing means the local error is not a valid
posterior interval inside the declared domain. It does not justify
clipping a Gaussian contour and calling it a new Fisher result. A bounded
inference calculation must use the same physical continuation and priors.

\paragraph{Checks needed before a scientific conclusion.}
Compare full derivatives at both step sizes and at refined response
sampling, with the covariance and physical choices held fixed. Then
compare the physical continuation alternatives. The mean uses the
numbered example's TATT choices, while the supplied covariance keeps its
own simulation assumptions. Quadratic IA currently retains the declared
gravity-only prescription. These are model assumptions to report, not
errors that vanish when the interpolation grid becomes denser. The
refactored notebook has no newly computed results; its eventual output
must state these conditions alongside any gain attributed to Map3.

\paragraph{Acknowledgements.}
\textbf{We thank Sunao Sugiyama and Rafael Gomes for their three-point codes
and analyses. We thank the BACCO authors and Pierre Burger and collaborators
for the feedback models and emulator work.} Source and model references are
\href{https://github.com/git-sunao/fastnc}{fastnc},
\href{https://github.com/rchgomes/3pcf_integrator}{3pcf\us integrator},
\href{https://github.com/rchgomes/desy6-3pcf}{desy6-3pcf}, and
\href{https://baccoemu.readthedocs.io/}{BACCO's documentation}.

\begin{thebibliography}{9}
\bibitem{methods} R. C. H. Gomes et al., \emph{Cosmology with second and third-order
shear statistics for the Dark Energy Survey: Methods and simulated analysis},
\href{https://arxiv.org/abs/2503.03964}{arXiv:2503.03964}.
\bibitem{burger} P. A. Burger et al., \emph{Euclid: An emulator for baryonic
effects on the matter bispectrum},
\href{https://arxiv.org/abs/2506.18974}{arXiv:2506.18974}.
\end{thebibliography}
\end{document}
